\documentclass[11pt]{article}
\usepackage[utf8]{inputenc}
\usepackage{amsmath}
\usepackage{amsfonts}
\usepackage{amssymb}
\usepackage{amsthm}
\usepackage{geometry}
\usepackage{hyperref}
\usepackage{xcolor}
\usepackage{algorithm}
\usepackage{algorithmic}
\usepackage{tikz}
\usepackage{pgfplots}

\geometry{margin=1in}
\hypersetup{
    colorlinks=true,
    linkcolor=blue,
    filecolor=magenta,      
    urlcolor=cyan,
}

\title{Mathematical Formulation of Attention-Guided Reinforcement Learning}
\author{AttentionGuidedRL Repository}
\date{\today}

\newtheorem{definition}{Definition}
\newtheorem{theorem}{Theorem}
\newtheorem{lemma}{Lemma}

\begin{document}

\maketitle

\begin{abstract}
This document provides a comprehensive mathematical formulation of the optimization objective implemented in the Attention-Guided Reinforcement Learning (AttentionGuidedRL) repository. The system trains a language model to autonomously guide its own training by sequencing key-value pairs from Wikipedia articles using reinforcement learning with policy gradients, multi-head attention mechanisms, and proximal policy optimization (PPO).
\end{abstract}

\tableofcontents

\section{Introduction}

The AttentionGuidedRL system implements a novel approach where a base language model (Llama-3.2-3B or GPT-2) learns to sequence its own training data through reinforcement learning. The model generates vector queries using attention mechanisms and selects key-value pairs to maximize learning progress, measured by conditional log probabilities of value tokens.

\section{Problem Formulation}

\subsection{State Space and Action Space}

\begin{definition}[State Space]
At timestep $t$, the state $s_t$ consists of:
\begin{align}
s_t = (c_t, \mathcal{K}_t^{\text{available}})
\end{align}
where:
\begin{itemize}
    \item $c_t \in \mathbb{Z}^{L_t}$ is the context sequence of length $L_t$
    \item $\mathcal{K}_t^{\text{available}} \subseteq \{1, 2, \ldots, K\}$ is the set of available key indices
\end{itemize}
\end{definition}

\begin{definition}[Action Space]
The action space $\mathcal{A}_t$ at timestep $t$ is the set of available key indices:
\begin{align}
\mathcal{A}_t = \mathcal{K}_t^{\text{available}}
\end{align}
\end{definition}

\subsection{Trajectory and Episode Structure}

A trajectory $\tau$ consists of $T$ steps:
\begin{align}
\tau = \{(s_1, a_1, r_1), (s_2, a_2, r_2), \ldots, (s_T, a_T, r_T)\}
\end{align}

where $T$ is the number of key-value pairs per episode (NUM\_KV\_PAIRS in the code).

\section{Policy Definition}

\subsection{Vector Query Generation}

The policy $\pi_\theta$ operates through a vector query mechanism parameterized by LoRA adapter weights $\theta$.

\begin{definition}[Query Vector Generation]
Given context $c_t$, the query vector is generated as:
\begin{align}
q_t = \text{AvgPool}(\text{Embed}_\theta(c_t \oplus [\text{``Query''}]))
\end{align}
where:
\begin{itemize}
    \item $\text{Embed}_\theta(\cdot)$ extracts embeddings from the second-to-last attention layer
    \item $\oplus$ denotes sequence concatenation
    \item $\text{AvgPool}$ averages over the sequence dimension
    \item $q_t \in \mathbb{R}^{d_{\text{model}}}$ where $d_{\text{model}}$ is the model dimension
\end{itemize}
\end{definition}

\subsection{Multi-Head Attention Similarity}

The system handles both standard Multi-Head Attention (MHA) and Grouped Query Attention (GQA).

\begin{definition}[Head-wise Similarity Computation]
For each attention head $h \in \{1, 2, \ldots, H\}$ and key $k$:
\begin{align}
\text{sim}_{t,k}^{(h)} = \frac{(q_t^{(h)})^T k_k^{(\text{group}(h))}}{\sqrt{d_{\text{head}}}} \cdot \frac{1}{\tau}
\end{align}
where:
\begin{itemize}
    \item $q_t^{(h)} \in \mathbb{R}^{d_{\text{head}}}$ is the query for head $h$
    \item $k_k^{(\text{group}(h))} \in \mathbb{R}^{d_{\text{head}}}$ is the key for the corresponding group
    \item $\text{group}(h) = \lfloor h \cdot G / H \rfloor$ maps heads to groups (for GQA)
    \item $G$ is the number of key-value groups
    \item $H$ is the number of query heads
    \item $d_{\text{head}} = d_{\text{model}} / H$ is the head dimension
    \item $\tau = 1.0$ is the temperature parameter
\end{itemize}
\end{definition}

\begin{definition}[Per-Head Probability Distribution]
For each head $h$, the probability distribution over keys is:
\begin{align}
p_{t,k}^{(h)} = \frac{\exp(\text{sim}_{t,k}^{(h)})}{\sum_{k' \in \mathcal{K}_t^{\text{available}}} \exp(\text{sim}_{t,k'}^{(h)})}
\end{align}
\end{definition}

\begin{definition}[Policy Distribution]
The final policy distribution is obtained by averaging over heads:
\begin{align}
\pi_\theta(k | s_t) = \frac{1}{H} \sum_{h=1}^{H} p_{t,k}^{(h)}
\end{align}
\end{definition}

\section{Reward Function}

\begin{definition}[Step Reward]
The reward at step $t$ is defined as the conditional log probability of value tokens:
\begin{align}
r_t = \begin{cases}
\log p_\theta(v_t | c_t, k_t) - \log p_{\text{ref}}(v_t | c_t, k_t) & \text{if SUBTRACT\_BASE\_MODEL\_LOGPROBS} \\
\log p_\theta(v_t | c_t, k_t) & \text{otherwise}
\end{cases}
\end{align}
where:
\begin{itemize}
    \item $v_t$ are the value tokens at step $t$
    \item $k_t$ are the selected key tokens at step $t$
    \item $p_\theta(\cdot)$ is the adapter model probability
    \item $p_{\text{ref}}(\cdot)$ is the reference model probability (base model without LoRA)
\end{itemize}
\end{definition}

\begin{definition}[Conditional Log Probability]
The conditional log probability is computed as:
\begin{align}
\log p_\theta(v_t | c_t, k_t) = \frac{1}{|v_t|} \sum_{i=1}^{|v_t|} \log p_\theta(v_{t,i} | c_t, k_t, v_{t,<i})
\end{align}
where $v_{t,i}$ is the $i$-th token in the value sequence and $v_{t,<i}$ represents the preceding tokens.
\end{definition}

\section{Optimization Objective}

\subsection{Primary Objective Function}

\begin{theorem}[Trajectory Sampling with Step-Level Credit Assignment]
The system maximizes the following objective function:
\begin{align}
\mathcal{J}(\theta) = \mathbb{E}_{\tau \sim \pi_\theta} \left[ \sum_{t=1}^T A_t(\tau) \cdot \log \pi_\theta(a_t | s_t) \right] - \beta \cdot D_{KL}(\pi_\theta || \pi_{\text{ref}})
\end{align}

where:
\begin{itemize}
    \item $\tau = (s_1, a_1, r_1, \ldots, s_T, a_T, r_T)$ is a complete trajectory sampled from policy $\pi_\theta$
    \item $A_t(\tau) = R_t(\tau) - b_t$ is the \textbf{trajectory-dependent advantage} at step $t$
    \item $R_t(\tau) = \sum_{t'=t}^T \gamma^{t'-t} r_{t'}$ are the discounted returns from step $t$ forward in trajectory $\tau$
    \item $b_t = \mathbb{E}_{\text{batch}}[R_t(\tau)]$ is the GRPO baseline (batch average of returns at step $t$)
    \item $\gamma = 0.99$ is the discount factor
    \item $\beta = 0.1$ is the KL penalty coefficient
    \item $\pi_{\text{ref}}$ is the fixed reference policy (base model without LoRA)
\end{itemize}

\textbf{Key Insight:} We sample complete trajectories first, then assign step-specific credit. Each action $a_t$ receives weight $A_t(\tau)$ that depends on the \textbf{actual future rewards} observed in that trajectory, not a trajectory-average.
\end{theorem}

\subsection{Return and Advantage Computation}

\begin{definition}[Discounted Returns]
The return at timestep $t$ is:
\begin{align}
R_t = \sum_{t'=t}^T \gamma^{t'-t} r_{t'}
\end{align}
\end{definition}

\begin{definition}[GRPO Advantage Estimation]
The system uses Group Relative Policy Optimization (GRPO) for advantage estimation:
\begin{align}
A_t(\tau^{(i)}) &= R_t(\tau^{(i)}) - b_t \\
b_t &= \frac{1}{B} \sum_{i=1}^B R_t(\tau^{(i)})
\end{align}
where:
\begin{itemize}
    \item $\tau^{(i)}$ is the $i$-th trajectory in the batch
    \item $b_t$ is the per-timestep baseline computed as the batch average of returns at step $t$
    \item $B$ is the batch size
    \item Each trajectory gets its own advantage based on its actual returns vs. the batch average
\end{itemize}
\end{definition}

\section{Loss Function and Training Algorithm}

\subsection{PPO Loss Implementation}

\begin{theorem}[Step-Level PPO Loss with Process Supervision]
The actual loss function implements step-level Proximal Policy Optimization where each action receives credit for its own contribution:
\begin{align}
\mathcal{L}(\theta) = -\sum_{t=1}^T \left[ \min\left( \rho_t(\theta) \cdot A_t, \text{clip}(\rho_t(\theta), 1-\epsilon, 1+\epsilon) \cdot A_t \right) \right] + \beta \cdot D_{KL}(\pi_\theta || \pi_{\text{ref}})
\end{align}
where each step $t$ contributes independently:
\begin{itemize}
    \item $A_t = R_t - b_t$ is the \textbf{step-specific advantage} (return from step $t$ forward minus baseline)
    \item $\rho_t(\theta) = \frac{\pi_\theta(a_t | s_t)}{\pi_{\text{old}}(a_t | s_t)}$ is the \textbf{step-specific probability ratio}
    \item $a_t$ is the key selection action at step $t$ given context $s_t$
    \item $\epsilon = 0.2$ is the PPO clipping parameter
    \item $\text{clip}(x, a, b) = \max(a, \min(x, b))$
    \item KL divergence is computed against the fixed reference policy $\pi_{\text{ref}}$
\end{itemize}

\textbf{Process Supervision Interpretation:} Each action $a_t$ (key selection) receives credit proportional to its step-specific advantage $A_t$, enabling proper credit assignment where early selections (affecting more future rewards) receive higher weight than later selections.
\end{theorem}

\subsection{Step-Level Credit Assignment Details}

\begin{definition}[Step-Level Gradient Decomposition]
The policy gradient decomposes naturally into step-level contributions:
\begin{align}
\nabla_\theta \mathcal{J}(\theta) = \sum_{t=1}^T \mathbb{E}\left[ A_t \cdot \nabla_\theta \log \pi_\theta(a_t | s_t) \right]
\end{align}
\end{definition}

This formulation ensures:
\begin{enumerate}
    \item \textbf{Proper Temporal Credit}: Action $a_1$ (first key selection) affects rewards $\{r_1, r_2, \ldots, r_T\}$, so $A_1 = \sum_{t'=1}^T \gamma^{t'-1} r_{t'} - b_1$
    \item \textbf{Diminishing Influence}: Action $a_T$ (last key selection) only affects $r_T$, so $A_T = r_T - b_T$
    \item \textbf{Independent Assessment}: Each action's contribution is evaluated against its own expected value, not averaged across the trajectory
\end{enumerate}

\textbf{Comparison with Trajectory-Level Formulation:}
\begin{itemize}
    \item \textbf{Trajectory-Level}: $\mathcal{L} = -\bar{A} \sum_{t=1}^T \log \pi_\theta(a_t | s_t)$ where $\bar{A} = \frac{1}{T}\sum_{t=1}^T A_t$
    \item \textbf{Step-Level}: $\mathcal{L} = -\sum_{t=1}^T A_t \log \pi_\theta(a_t | s_t)$ (our formulation)
\end{itemize}

The step-level formulation provides \textbf{finer-grained supervision} where early actions naturally receive higher gradients due to their broader impact on trajectory outcomes.

\subsection{KL Divergence Regularization}

\begin{definition}[KL Divergence Term]
The KL divergence regularization is computed as:
\begin{align}
D_{KL}(\pi_\theta || \pi_{\text{ref}}) = \sum_{t=1}^T \sum_{k \in \mathcal{K}_t^{\text{available}}} \pi_\theta(k|s_t) \log \frac{\pi_\theta(k|s_t)}{\pi_{\text{ref}}(k|s_t)}
\end{align}
\end{definition}

This is implemented using PyTorch's \texttt{F.kl\_div} with \texttt{log\_target=True} for numerical stability. The reference policy $\pi_{\text{ref}}$ is the fixed base model without LoRA, providing a stable regularization target.

\section{Training Algorithm}

\begin{algorithm}
\caption{Attention-Guided RL Training}
\begin{algorithmic}[1]
\REQUIRE Base model, tokenizer, Wikipedia dataset
\REQUIRE Hyperparameters: $\gamma, \beta, \epsilon$, learning rate, batch size
\STATE Initialize LoRA adapter parameters $\theta$
\STATE Initialize reference model $\pi_{\text{ref}}$ (base model without LoRA, fixed)
\STATE Initialize old model $\pi_{\text{old}} \leftarrow \pi_\theta$
\STATE Initialize optimizer (AdamW with learning rate $5 \times 10^{-4}$)
\FOR{episode $= 1$ to NUM\_EPISODES}
    \STATE Sample Wikipedia article and extract key-value pairs
    \STATE Generate trajectory $\tau$ using current policy $\pi_\theta$
    \STATE Compute rewards $\{r_t\}_{t=1}^T$ using Equation (7)
    \STATE Compute returns $\{R_t\}_{t=1}^T$ using Equation (10)
    \STATE Compute advantages $\{A_t\}_{t=1}^T$ using GRPO (Equation 11-12)
    \STATE Compute PPO loss $\mathcal{L}(\theta)$ using Equation (13) with KL vs $\pi_{\text{ref}}$
    \STATE Update parameters: $\theta \leftarrow \theta - \nabla_\theta \mathcal{L}(\theta)$
    \IF{episode $\bmod$ BASELINE\_UPDATE\_FREQUENCY $== 0$}
        \STATE Update old model: $\pi_{\text{old}} \leftarrow \pi_\theta$ (for PPO ratios only)
    \ENDIF
\ENDFOR
\end{algorithmic}
\end{algorithm}

\section{Implementation Details}

\subsection{Model Architecture}
\begin{itemize}
    \item \textbf{Base Models}: Llama-3.2-3B or GPT-2
    \item \textbf{LoRA Configuration}: rank=8, $\alpha$=16, dropout=0.0
    \item \textbf{Precision}: bfloat16 for computation, float32 for similarity calculations
    \item \textbf{Context Window}: Up to model maximum (2048 for GPT-2, 4096+ for Llama)
\end{itemize}

\subsection{Training Configuration}
\begin{itemize}
    \item \textbf{Optimizer}: AdamW with learning rate $5 \times 10^{-4}$
    \item \textbf{Batch Size}: 4 trajectories per update
    \item \textbf{Gradient Clipping}: Norm clipped to 1.0
    \item \textbf{Episodes}: 10,000 total training episodes
    \item \textbf{Checkpoint Frequency}: Every 100 episodes
\end{itemize}

\subsection{Data Processing}
\begin{itemize}
    \item \textbf{Dataset}: Wikipedia 2022-03-01 English subset
    \item \textbf{Token Counts}: 10 tokens per key, 10 tokens per value
    \item \textbf{Trajectory Length}: Up to 10 key-value pairs per episode
    \item \textbf{Prefixes}: "Key: " and "Value: " for context building
\end{itemize}

\section{Key Mathematical Properties}

\subsection{Convergence Properties}
\begin{lemma}[Policy Improvement]
Under the PPO clipped objective with appropriate clipping parameter $\epsilon$, the policy improvement is guaranteed to be conservative, preventing destructively large updates.
\end{lemma}

\subsection{Attention Mechanism Properties}
\begin{lemma}[Multi-Head Averaging]
The averaging operation over attention heads in Equation (6) preserves the probability distribution property:
\begin{align}
\sum_{k \in \mathcal{K}_t^{\text{available}}} \pi_\theta(k | s_t) = 1
\end{align}
\end{lemma}

\subsection{GRPO Baseline Properties}
\begin{theorem}[Unbiased Advantage Estimation]
The GRPO baseline provides unbiased gradient estimates:
\begin{align}
\mathbb{E}[A_t] = 0 \quad \text{for each timestep } t
\end{align}
This property holds by construction since $A_t = R_t - \mathbb{E}[R_t]$ where the expectation is over the batch.
\end{theorem}

\section{Conclusion}

This mathematical formulation describes a sophisticated reinforcement learning system that combines:
\begin{enumerate}
    \item \textbf{Vector-based policy}: Using attention layer embeddings for query generation
    \item \textbf{Multi-head attention similarity}: Handling both MHA and GQA architectures
    \item \textbf{Step-level PPO optimization}: With proper temporal credit assignment through step-specific advantages
    \item \textbf{GRPO advantage estimation}: For variance reduction without additional parameters
    \item \textbf{Process supervision}: Where each action receives credit proportional to its impact on future outcomes
    \item \textbf{Self-supervised learning}: Where the model learns to sequence its own training data
\end{enumerate}

The key innovation is the \textbf{step-level formulation} that ensures proper credit assignment: early key selections (which affect more future rewards) naturally receive higher gradient magnitudes than later selections. This creates a natural curriculum where the model learns to prioritize the quality of early decisions while still optimizing the entire sequence.

The system represents a novel approach to active learning where the model's attention mechanisms directly drive the selection of training examples, creating a feedback loop between the model's internal representations and its learning curriculum.

\bibliographystyle{plain}
\begin{thebibliography}{9}

\bibitem{schulman2017proximal}
Schulman, J., Wolski, F., Dhariwal, P., Radford, A., \& Klimov, O. (2017).
\textit{Proximal policy optimization algorithms}.
arXiv preprint arXiv:1707.06347.

\bibitem{shazeer2019fast}
Shazeer, N. (2019).
\textit{Fast transformer decoding: One write-head is all you need}.
arXiv preprint arXiv:1911.02150.

\bibitem{hu2021lora}
Hu, E. J., Shen, Y., Wallis, P., Allen-Zhu, Z., Li, Y., Wang, S., ... \& Chen, W. (2021).
\textit{LoRA: Low-rank adaptation of large language models}.
arXiv preprint arXiv:2106.09685.

\end{thebibliography}

\end{document} 